The first two arguments below use ordinary Tarskian signatures and
fragments. Their auxiliary signatures need not be objects of the
language category \(\mathcal C\). The final argument returns to the
category and its distinguished universe predicate explicitly.

\begin{annexitem}{annex:p03.ib.1}{thm:p03-ib-rank}{A bound from small pointed models}
\begin{proof}
Use fragment elementarity as in
\cref{def:p03-w03-p03-039-fragment}.
\Cref{thm:p03-w03-p03-039-downward-ls} gives an
\(F\)-elementary submodel of size at most \(\lambda\)
containing any prescribed set of at most \(\lambda\) elements.
\Cref{thm:p03-w03-p03-039-elementary-chain} preserves fragment
truth in unions of elementary chains. The finiteness of free-variable
lists is essential to both uses: every parameter tuple lies in one
stage, even when a formula has an infinite conjunction.

Let \(S\) be a set of representatives of the isomorphism types of
pairs \((\mathfrak M,a)\), where \(\mathfrak M\models\psi\),
\(|M|\leq\lambda\), and \(a\in D^{\mathfrak M}\).
This is a set of size at most \(2^\lambda\). Indeed every such
structure can be copied onto a subset of \(\lambda\). Each
finitary relation or function graph is then a subset of a finite
power of \(\lambda\), and there are at most \(\lambda\) symbols.
There are at most \((2^\lambda)^\lambda=2^\lambda\) possible
interpretations, and marking one point does not increase this
bound. Take representatives among these coded structures.

Put \((\mathfrak N,b)\triangleleft(\mathfrak M,a)\) when an
\(F\)-elementary embedding \(f:\mathfrak M\to\mathfrak N\)
exists such that
\[
 \mathfrak N\models\rho\bigl(b,f(a)\bigr).
\]
The relation does not depend on the representatives chosen:
isomorphisms transport the embedding and its distinguished points.
Because \(\delta\in F\), the image \(f(a)\) is again a
carrier point.

We claim that \(\triangleleft\) is well-founded. Otherwise,
Choice gives a sequence \((\mathfrak M_n,a_n)\) and witnessing
embeddings \(f_n:\mathfrak M_n\to\mathfrak M_{n+1}\) with
\(a_{n+1}\) below \(f_n(a_n)\). Replace the structures successively
by isomorphic copies in which these embeddings become inclusions.
At each step identify the image of the preceding structure with
that structure, and use fresh elements for the remaining points;
transport all interpretations along this bijection.
The resulting chain is \(F\)-elementary. Its union is a model
of \(\psi\) by the elementary-chain lemma. In that union the
images of the \(a_n\) form a strictly descending sequence in the
interpreted order: both the carrier predicate \(\delta\) and
comparison formula \(\rho\) are preserved. This contradicts
the assumed well-ordering. Thus the relation is well-founded.

Let \(r(s)\) be its well-founded rank, defined by
\[
 r(s)=\sup\{r(t)+1:t\triangleleft s\}.
\]
Set \(\Theta=(2^\lambda)^+\). This successor cardinal is regular.
Well-founded induction gives \(r(s)<\Theta\) for every \(s\),
since each predecessor family has size at most \(2^\lambda\).
The same cardinality bound on \(S\) gives
\[
 \eta=\sup_{s\in S}(r(s)+1)<\Theta.
\]
The empty set of types gives \(\eta=0\). As \(\Theta\) is an
infinite cardinal, \(\eta+1<\Theta\).

Fix now an arbitrary model \(\mathfrak A\models\psi\), with
no bound on its size. For \(a\in D^{\mathfrak A}\), let
\(o_{\mathfrak A}(a)\) be the order type of its predecessors.
We prove, by induction in this well-order, that every small
\(F\)-elementary \(\mathfrak M\prec_F\mathfrak A\) containing
\(a\) satisfies
\[
 r([\mathfrak M,a])\geq o_{\mathfrak A}(a).
\]
Brackets denote its representative type in \(S\).
For each \(b\) below \(a\), the fragment L\"owenheim--Skolem
theorem gives \(\mathfrak N\prec_F\mathfrak A\) of size at most
\(\lambda\) containing \(M\cup\{b\}\).
The inclusion \(\mathfrak M\subseteq\mathfrak N\) is
\(F\)-elementary, since both structures agree with
\(\mathfrak A\) on all fragment formulas with parameters in
\(M\). Hence
\([\mathfrak N,b]\triangleleft[\mathfrak M,a]\).
The induction hypothesis and the definition of rank yield
\[
\begin{aligned}
 r([\mathfrak M,a])
   &\geq\sup_{b\,\rho^{\mathfrak A}\,a}
                     \bigl(o_{\mathfrak A}(b)+1\bigr)\\
   &=o_{\mathfrak A}(a).
\end{aligned}
\]
Here the induction hypothesis applies to every small submodel
containing the lower point, so the newly chosen \(\mathfrak N\)
is covered. There is no requirement that \(M\) already contain
all predecessors of \(a\).

For each carrier point choose such a small \(\mathfrak M\).
Its rank is among those used to define \(\eta\), so
\(o_{\mathfrak A}(a)+1\leq\eta\). Taking the supremum over
all carrier points shows that the interpreted order type is at
most \(\eta\). An empty carrier has type zero and satisfies
the same bound. Thus \(\beta=\eta+1\) proves the assertion.
\end{proof}
\end{annexitem}

\begin{annexitem}{annex:p03.ib.2}{thm:p03-ib-individual}{Individual bounds and an unbounded supremum}
\begin{proof}
Fix a sentence \(\varphi\) and its designated binary relation
\(<\). If the distinguished universe symbol \(\mathbf V\) is
present, conjoin \(\forall x\,\mathbf V(x)\) when treating it
in an ordinary signature; this leaves its intended models unchanged.
Retain only the symbols actually occurring in \(\varphi\)
and \(<\). Let
\(\delta(x)=\exists y\,(x<y\lor y<x)\), the field formula.
Choose an infinite cardinal \(\lambda\) bounding this vocabulary
and the set of subformulas of \(\varphi\). By
\cref{thm:p03-w03-p03-039-countable-fragment-size}, the sentence
and field and order formulas lie in a fragment of size at most
\(\lambda\), with finitely many free variables in each formula.
Every model of \(\varphi\) has the required interpreted
well-order, so \cref{thm:p03-ib-rank} supplies a uniform bound.
Unused vocabulary symbols do not change satisfaction. Nonempty
structures for the retained vocabulary can be expanded to the
original vocabulary by interpreting additional functions constantly
at one chosen point and additional relations arbitrarily.

For the other assertion extend the ordinal-formula construction
of \cref{def:p03-w03-p03-041-mu} to all ordinals. Simultaneously
construct the two variable-renamed versions of
\[
 \mu_\beta(x)=\forall y
       \left(y<x\leftrightarrow
                   \bigvee_{\gamma<\beta}\mu_\gamma(y)\right).
\]
When changing the displayed free variable, swap the two variable
names throughout, including their bound occurrences. Each formula
has only its displayed variable free. Transfinite recursion uses
sets of earlier formulas and well-founded trees, so these are
legitimate \(L_{\infty\omega}\)-formulas, not proper-class
conjunctions.

In a strict linear order, transfinite induction proves that
\(\mu_\beta(a)\) holds exactly when the predecessor order of
\(a\) has type \(\beta\). For completeness, if that order has
type \(\beta\), every predecessor has some smaller rank and
no non-predecessor has a smaller predecessor rank, giving the
biconditional. Conversely, the biconditional labels every
predecessor by a rank below \(\beta\), using the induction
hypothesis. These labels preserve and reflect the order, since
initial segments of a linear order are nested. Their image is
an initial segment of \(\beta\), so the predecessor order has
some type \(\zeta\leq\beta\). If \(\zeta<\beta\), the
induction hypothesis gives \(\mu_\zeta(a)\), making the
right side of the biconditional true at \(y=a\), contrary to
irreflexivity. Hence \(\zeta=\beta\).

For \(\alpha\geq2\), conjoin the strict linear-order axioms
on the universe with
\[
 \forall x\,\bigvee_{\beta<\alpha}\mu_\beta(x).
\]
Every model is a well-order of type at most \(\alpha\),
by the predecessor-rank description. The ordinal \(\alpha\)
itself is a model, and its strict order has its whole universe
as field because \(\alpha\geq2\). Thus this sentence pins
\(\alpha\). As \(\alpha\) is arbitrary, the supremum of the
pinnable ordinals is \(\infty\), despite the individual bounds.
\end{proof}
\end{annexitem}

\begin{annexitem}{annex:p03.ib.3}{thm:p03-w03-p03-048-lopez-escobar}{The Lopez-Escobar bound}
\begin{proof}
Let a class of ordinals be \(\Sigma^1_1\) in
\(\mathcal L_{\lambda\omega}\) in the sense of
\cref{def:p03-w03-p03-048-sigma,def:p03-w03-p03-048-infinitary}.
It is witnessed by a sentence \(\psi\) over a language \(K\)
and a zero-term pullback from the pure-order language
\(L_< =\{\mathbf V,<\}\). Write \(P\) for the unary image
of \(\mathbf V\), and \(R\) for the binary image of \(<\).
The interpreted order is exactly \(R\) restricted to \(P\).
There are no function symbols in the pure-order language, so
its pullback requires no additional closure condition. The
empty interpreted universe is allowed and gives order type zero.

Let \(K_0\) consist of the symbols in \(\psi\) and
\(\mathbf V,P,R\). Its size is at most
\(\lambda\), since \(\psi\) has fewer than \(\lambda\)
subformulas and every atomic formula contains only finitely
many symbols. Put
\(\psi_0=\psi\land\forall x\,\mathbf V(x)\), and regard this
as an ordinary Tarskian formula over the sub-signature. Every
nonempty \(K_0\)-model of \(\psi_0\)
expands to a \(K\)-model: choose one element for new constants
and constant-valued new functions, and arbitrary new relations.
The distinguished universe predicate remains the whole universe.
Tarskian reduct invariance preserves \(\psi\), and the
interpretations of \(P,R\) are unchanged. Hence every such
reduced model interprets a well-order, by the hypothesis that
the pullback class consists only of ordinals.
An empty reduced model, if allowed, has empty carrier and
contributes only order type zero.

This auxiliary sub-signature need not be an object of
\(\mathcal C\). It is used only for the ordinary Tarskian
fragment construction, and the expansion just described returns
to the original \(K\)-structures. Thus no extra closure condition
on the source's language category is assumed.

By \cref{thm:p03-w03-p03-039-countable-fragment-size}, choose
an infinitary fragment \(F\) of size at most \(\lambda\)
containing \(\psi_0\), \(P(x)\), \(R(x,y)\), and the atomic
formulas of this sub-signature. Use
\(\delta(x)=P(x)\) and \(\rho(x,y)=R(x,y)\) in
\cref{thm:p03-ib-rank}. All interpreted order types are below
some \(\beta<(2^\lambda)^+\), and therefore below
\((2^\lambda)^+\). The original class is among these types,
so it has the asserted bound. This works for every such class,
giving \(\operatorname{wo}(\mathcal L_{\lambda\omega})
\leq(2^\lambda)^+\) with the strict memberwise meaning of
boundedness fixed in the chapter.
\end{proof}
\end{annexitem}
